% Глава: 14-make
% Статус: наполнено (первый прогон).
\chapter{Система сборки: make}
\label{ch:14-make}

\section{Зачем нужен make}

Когда проект вырастает за один \cmd{gcc}, ручная сборка
становится ненадёжной: легко забыть файл или пересобрать
то, что не менялось. \textbf{make} --- система сборки:
он смотрит на \emph{зависимости} (что из чего) и
пересобирает \textbf{только изменившееся}.

Выполняем сборку на своей ВМ: утилита \cmd{make} читает
\fpath{Makefile} на сервере \fpath{pglab.nsu.ru} и пересобирает нужное.

\section{Анатомия Makefile}

Файл \fpath{Makefile} состоит из \textbf{правил}
(rules):
\begin{minted}[bgcolor=gray!6]{Makefile}
target: prerequisites
    recipe command
\end{minted}

\begin{itemize}
  \item \textbf{target} --- что строим (файл или имя);
  \item \textbf{prerequisites} --- от чего зависит;
  \item \textbf{recipe} --- команды (всегда с \emph{табуляцией},
    не с пробелами!).
\end{itemize}

\begin{warnbox}
В листингах этой книги отступ рецепта показан \textbf{четырьмя
пробелами} --- табуляцию внутри блока напечатать нельзя. В настоящем
\file{Makefile} перед командой рецепта \textbf{обязательно стоит
одна таб}: make с пробелами в начале строки не примет и выдаст
ошибку \emph{missing separator}.
\end{warnbox}

Минимальный Makefile для \fpath{hello.c}:
\begin{minted}[bgcolor=gray!6]{Makefile}
CC = gcc
CFLAGS = -Wall -Wextra -g

hello: hello.c
    $(CC) $(CFLAGS) -o hello hello.c

clean:
    rm -f hello
\end{minted}

\begin{warnbox}
Рецепты в Makefile \textbf{обязаны} начинаться с табуляции.
Пробелы --- синтаксическая ошибка. При копировании из
редактора проверяйте, что это именно табы.
\end{warnbox}

\section{Переменные и автоматические переменные}

Обычные переменные задаются как \cmd{CC = gcc}.
Автоматические переменные зависят от текущего правила:
\begin{center}
\begin{tabular}{lp{0.6\linewidth}}
\toprule
\textbf{Переменная} & \textbf{Значение} \\
\midrule
\cmd{$@}   & целевой файл текущего правила \\
\cmd{$<}   & первый из зависимостей \\
\cmd{$^}   & все зависимости \\
\cmd{$(CC)} & значение переменной \\
\bottomrule
\end{tabular}
\end{center}

Использование \cmd{\$@} и \cmd{\$<} укорачивает правила:
\begin{minted}[bgcolor=gray!6]{Makefile}
app: main.o util.o
    $(CC) -o $@ $^
\end{minted}

\section{Паттерн-правила}

Паттерн-правило описывает, как строить \emph{любой} файл
по шаблону. Классический пример: из \fpath{.c} в \fpath{.o}:
\begin{minted}[bgcolor=gray!6]{Makefile}
%.o: %.c
    $(CC) $(CFLAGS) -c $< -o $@
\end{minted}

\cmd{\%} --- заглушка, подставляющая имя.

\section{Типовой Makefile для проекта}

\begin{minted}[bgcolor=gray!6]{Makefile}
CC      = gcc
CFLAGS  = -Wall -Wextra -g
SRCDIR  = src
OBJDIR  = obj
SRC     = $(wildcard $(SRCDIR)/*.c)
OBJ     = $(SRC:$(SRCDIR)/%.c=$(OBJDIR)/%.o)
DEPS    = $(OBJ:.o=.d)

all: app

app: $(OBJ)
    $(CC) $(CFLAGS) -o $@ $^

$(OBJDIR)/%.o: $(SRCDIR)/%.c | $(OBJDIR)
    $(CC) $(CFLAGS) -MMD -MP -c $< -o $@

$(OBJDIR):
    mkdir -p $(OBJDIR)

-include $(DEPS)

clean:
    rm -rf $(OBJDIR) app

.PHONY: all clean
\end{minted}

Обратите внимание: \cmd{.PHONY} говорит make, что
\cmd{all} и \cmd{clean} --- не файлы; это важно, если
на диске появятся файлы с такими именами.

\section{Полезные флаги \cmd{make}}

\begin{minted}[bgcolor=gray!6]{bash}
make                 # целевая по умолчанию, обычно all
make clean           # выполнить целевую clean
make -n              # вывести команды без выполнения (dry run)
make -j4             # запускать до 4 задач параллельно
make -f other.mk     # использовать другой makefile
make CC=clang        # переопределить переменную в командной строке
\end{minted}

\section{Типичные ошибки}
\begin{itemize}
  \item \emph{missing separator} --- в рецепте пробелы
    вместо табуляции;
  \item Циклическая зависимость --- make останавливается
    с ошибкой;
  \item Пересборка каждый раз --- пропущена зависимость
    (например, заголовочный файл); используйте
    \cmd{-MMD -MP} для автоматического отслеживания;
  \item \cmd{make clean} делает что-то лишнее --- проверьте,
    что в рецепте только нужные команды.
\end{itemize}

\section{Связь с остальным курсом}

\cmd{make} --- стандарт сборки C-проектов
(\nameref{ch:13-c-compilation}) и в Docker-образах
(\nameref{ch:20-dockerfile}). Понимание зависимостей
пригодится и в скриптах CI
(\nameref{ch:26-ci}).

\section{Задачи для самопроверки}
\begin{enumerate}
  \item Из чего состоит правило в Makefile?
  \item Зачем \cmd{.PHONY}?
  \item Чем \cmd{\$@} отличается от \cmd{\$<}?
  \item Что делает \cmd{make -n}?
  \item Как заставить make пересобирать проект при
    изменении заголовочного файла?
\end{enumerate}
